% Part: first-order-logic
% Chapter: completeness
% Section: construction-of-model

\documentclass[../../../include/open-logic-section]{subfiles}

\begin{document}

\iftag{FOL}
      {\olfileid{fol}{com}{mod}}
      {\olfileid{pl}{com}{mod}}

\olsection{Een model bouwen}

\begin{explain}
\iftag{FOL}{Voorlopig laten we~$\eq$ buiten beschouwing. We willen dus
  alleen bewijzen dat een consistente verzameling~$\Gamma$ van
  !!{sentence}s zonder~$\eq$ vervulbaar is. We breiden~$\Gamma$ uit tot
  een verzameling~$\Gamma^*$ die consistent, !!{complete} en verzadigd is.
  Daaruit bouwen we een model~$\Struct{M(\Gamma^*)}$. Als domein nemen we
  alle gesloten termen van~$\Lang{L'}$. Elke !!{constant} verwijst naar
  zichzelf. Sterker nog: voor elke gesloten term~$t$ geldt
  $\Value{t}{M(\Gamma^*)} = t$. Aan elk !!{predicate} geven we zo'n
  extensie dat een atomaire !!{sentence} precies dan waar is in
  $\Struct{M(\Gamma^*)}$ wanneer zij in~$\Gamma^*$ zit. Daardoor zijn alle
  atomaire !!{sentence}s uit~$\Gamma^*$ waar in
  $\Struct{M(\Gamma^*)}$. De andere
  zinnen zijn ook waar, omdat~$\Gamma^*$ consistent, volledig en
  verzadigd is.}{We kunnen nu !!a{valuation} maken die elk
  $!A \in \Gamma$ waar maakt. Eerst gebruiken we het lemma van Lindenbaum.
  Dat geeft een volledige consistente verzameling
  $\Gamma^* \supseteq \Gamma$. Daarna laten we de
  !!{propositional variable}s in~$\Gamma^*$ bepalen welke
  waardering~$\pAssign v(\Gamma^*)$ we nemen.}
\end{explain}

\iftag{FOL}{%
\begin{defn}[Termmodel]
\ollabel{defn:termmodel}
Neem een !!{complete}, consistente en verzadigde verzameling
$\Gamma^*$ van !!{sentence}s in een taal~$\Lang L$. Het
\emph{termmodel}~$\Struct M(\Gamma^*)$ van $\Gamma^*$ is de volgende
!!{structure}:
\begin{enumerate}
\item Het !!{domain}~$\Domain{M(\Gamma^*)}$ bestaat uit alle gesloten
  termen van~$\Lang L$.
\item Een !!{constant} $c$ verwijst naar $c$ zelf:
  $\Assign{c}{M(\Gamma^*)} = c$.
\item Krijgt de !!{function}~$f$ de gesloten termen $t_1$, \dots, $t_n$
  als invoer, dan is haar uitkomst de gesloten term
  $f(t_1, \dots, t_n)$:
\[
\Assign{f}{M(\Gamma^*)}(t_1, \dots, t_n) = f(t_1,\dots, t_n)
\]
\item Is $R$ een $n$-plaatsig !!{predicate}, dan geldt
  \[
  \tuple{t_1, \dots,
  t_n} \in \Assign{R}{M(\Gamma^*)} \text{ precies als } \Atom{R}{t_1, \dots,
    t_n} \in \Gamma^*.
  \]
\end{enumerate}
\end{defn}
}{%
\begin{defn}
\ollabel{defn:termmodel}
Neem een volledige consistente verzameling $\Gamma^*$ van
!!{formula}s. We definiëren
\[
\pAssign v(\Gamma^*)(p) = \begin{cases}
  \True & \text{als $p \in \Gamma^*$}\\
  \False & \text{als $p \notin \Gamma^*$}
  \end{cases}
\]
\end{defn}
}

\iftag{FOL}{%
Nu controleren we dat werkelijk $\Value{t}{M(\Gamma^*)} = t$.

\begin{lem} 
\ollabel{lem:val-in-termmodel} Neem het termmodel
$\Struct M(\Gamma^*)$ uit \olref{defn:termmodel}. Dan geldt
$\Value{t}{M(\Gamma^*)} = t$.
\end{lem}
\begin{proof} 
 We bewijzen dit met inductie naar $t$. Eerst het basisgeval: als $t$
 !!a{constant} is, volgt de bewering meteen uit de definitie van het
 termmodel. Neem voor de volgende stap gesloten termen
 $t_1, \ldots, t_n$ waarvoor $\Value{t_i}{M(\Gamma^*)} = t_i$, en neem
 een !!{function}~$f$ met $n$ argumenten. Dan
\begin{align*}
\Value{f(t_1,\ldots,t_n)}{M(\Gamma^*)} &= \Assign{f}{M(\Gamma^*)}(\Value{t_1}
 {M(\Gamma^*)}, 
\ldots, \Value{t_n}{M(\Gamma^*)}) \\
&= \Assign{f}{M(\Gamma^*)}(t_1, \dots, t_n) \\
&= f(t_1,\dots, t_n),
\end{align*} 
De inductie laat dus zien dat dit voor elke gesloten term~$t$ geldt.
\end{proof}
}

\iftag{FOL}{%
\begin{explain}
  \iftag{prvEx}{Een !!{structure}~$\Struct{M}$ kan een existentieel
    gekwantificeerde !!{sentence}~$\lexists[x][!A(x)]$ waar maken zonder
    ook maar één instantie~$!A(t)$ waar te maken.}{}
  \iftag{prvAll}{Een !!{structure}~$\Struct{M}$ kan elke instantie
    $!A(t)$ van een universeel gekwantificeerde
    !!{sentence}~$\lforall[x][!A(x)]$ waar maken, terwijl
    $\lforall[x][!A(x)]$ zelf niet waar is.}{} De reden is dat niet elk
  !!{element} van~$\Domain{M}$ de waarde van een gesloten term hoeft te
  zijn: $\Struct{M}$ hoeft niet overdekt te zijn. Daarom gebruikt de
  definitie van vervulling toekenningen van waarden aan variabelen. Bij
  ons termmodel~$\Struct{M(\Gamma^*)}$ is dat probleem er niet. Dit model
  is namelijk wel overdekt. Het volgende resultaat zegt precies wat
  daaruit volgt.
\end{explain}

\begin{prop}
\ollabel{prop:quant-termmodel}
Neem het termmodel $\Struct M(\Gamma^*)$ uit \olref{defn:termmodel}.
\begin{tagenumerate}{prvEx,prvAll}
\tagitem{prvEx}{$\Sat{M(\Gamma^*)}{\lexists[x][!A(x)]}$ precies als
  $\Sat{M(\Gamma^*)}{!A(t)}$ voor minstens één gesloten term~$t$.}{}
\tagitem{prvAll}{$\Sat{M(\Gamma^*)}{\lforall[x][!A(x)]}$ precies als
  $\Sat{M(\Gamma^*)}{!A(t)}$ voor elke gesloten term~$t$.}{}
\end{tagenumerate}
\end{prop}

\begin{proof}
\begin{tagenumerate}{prvEx,prvAll}
\tagitem{prvEx}{%
  \iftag{probEx}{Opgave.}{Volgens \olref[syn][ass]{prop:sat-quant} geldt
    $\Sat{M(\Gamma^*)}{\lexists[x][!A(x)]}$ precies als er minstens één
    toekenning~$s$ van waarden aan variabelen is waarvoor
    $\Sat{M(\Gamma^*)}{!A(x)}[s]$. Het domein $\Domain{M(\Gamma^*)}$
    bestaat uit de gesloten termen van~$\Lang{L}$. De vorige voorwaarde
    geldt daarom precies als er minstens één gesloten term~$t$ is met
    $s(x) = t$ en $\Sat{M(\Gamma^*)}{!A(x)}[s]$. Volgens
    \olref[fol][syn][ext]{prop:ext-formulas},
    geldt $\Sat{M(\Gamma^*)}{!A(x)}[s]$ precies als
    $\Sat{M(\Gamma^*)}{!A(t)}[s]$, waarbij $s(x) = t$. Volgens
    \olref[fol][syn][ass]{prop:sentence-sat-true} geldt
    $\Sat{M(\Gamma^*)}{!A(t)}[s]$ precies als
    $\Sat{M(\Gamma^*)}{!A(t)}$, want $!A(t)$ is een zin.}}{}
\tagitem{prvAll}{%
  \iftag{probAll}{Opgave.}{Volgens \olref[syn][ass]{prop:sat-quant} geldt
    $\Sat{M(\Gamma^*)}{\lforall[x][!A(x)]}$ precies als voor elke
    toekenning $s$ van waarden aan variabelen,
    $\Sat{M(\Gamma^*)}{!A(x)}[s]$. Het domein $\Domain{M(\Gamma^*)}$
    bestaat uit de gesloten termen van~$\Lang{L}$. Bij elke gesloten
    term~$t$ kunnen we dus $s(x) = t$ kiezen. Omgekeerd is $s(x)$ bij
    elke toekenning een bepaalde gesloten term~$t$. Volgens
    \olref[fol][syn][ext]{prop:ext-formulas} geldt
    $\Sat{M(\Gamma^*)}{!A(x)}[s]$ precies als
    $\Sat{M(\Gamma^*)}{!A(t)}[s]$, waarbij $s(x) = t$. Volgens
    \olref[fol][syn][ass]{prop:sentence-sat-true} geldt
    $\Sat{M(\Gamma^*)}{!A(t)}[s]$ precies als
    $\Sat{M(\Gamma^*)}{!A(t)}$, want $!A(t)$ is een zin.}}{}
\end{tagenumerate}
\end{proof}
}{}

\tagprob[FOL]{probEx,probAll}
\begin{prob}
Maak het bewijs van \olref[fol][com][mod]{prop:quant-termmodel} af.
\end{prob}
\tagendprob

\begin{lem}[Waarheidslemma]
\ollabel{lem:truth} \iftag{FOL}{Neem aan dat $!A$ geen~$\eq$ bevat. Dan}{}
geldt $\iftag{FOL}{\Sat{M(\Gamma^*)}}{\pSat{v(\Gamma^*)}}{!A}$ precies als
$!A \in \Gamma^*$.
\end{lem}

\begin{proof}
We bewijzen de twee richtingen tegelijk, met inductie naar $!A$.
\begin{enumerate}
\tagitem{prvFalse}{\indcase{!A}{\lfalse}
  {$\iftag{FOL}{\Sat/{M(\Gamma^*)}{\lfalse}}{\pSat/{v(\Gamma^*)}{\lfalse}}$
    volgens de definitie van vervulling. Tegelijk is
    $\lfalse \notin \Gamma^*$, want $\Gamma^*$ is consistent.}}{}

\tagitem{prvTrue}{\indcase{!A}{\ltrue}
  {$\iftag{FOL}{\Sat{M(\Gamma^*)}}{\pSat{v(\Gamma^*)}}{\ltrue}$
    volgens de definitie van vervulling. Ook is $\ltrue \in \Gamma^*$:
    de verzameling $\Gamma^*$ is consistent en !!{complete}, en
    $\Gamma^* \Proves \ltrue$.}}{}

\tagitem{FOL}{\indcase{!A}{R(t_1, \dots, t_n)}
  {$\Sat{M(\Gamma^*)}{\Atom{R}{t_1, \dots, t_n}}$ precies als
    $\tuple{t_1, \dots, t_n} \in \Assign{R}{M(\Gamma^*)}$, volgens de
    definitie van vervulling. Dat geldt weer precies als
    $R(t_1, \dots, t_n) \in \Gamma^*$, volgens de bouw van
    $\Struct M(\Gamma^*)$.}}{
    \indcase{!A}{p}{$\pSat{v(\Gamma^*)}{p}$ precies als
    $\pAssign v(\Gamma^*)(p) = \True$, volgens de definitie van
    vervulling. Dat geldt weer precies als $p \in \Gamma^*$, volgens de
    bouw van $\pAssign v(\Gamma^*)$.}}

\tagitem{prvNot}{%
  \indcase{!A}{\lnot !B}
    {$\iftag{FOL}{\Sat{M(\Gamma^*)}}{\pSat{v(\Gamma^*)}}{\indfrm}$ precies
    als $\iftag{FOL}{\Sat/{M(\Gamma^*)}{!B}}{\pSat/{v(\Gamma^*)}{!B}}$,
    volgens de definitie van vervulling. Volgens de inductieaanname geldt
    $\iftag{FOL}{\Sat/{M(\Gamma^*)}{!B}}{\pSat/{v(\Gamma^*)}{!B}}$ precies
    als $!B \notin \Gamma^*$. Omdat $\Gamma^*$ consistent en
    !!{complete} is, geldt $!B \notin \Gamma^*$ precies als
    $\lnot !B \in \Gamma^*$.}}{}

\tagitem{prvAnd}{%
\iftag{probAnd}{%
  \indcase!{!A}{!B \land !C}{}}{%
  \indcase{!A}{!B \land
    !C}{$\iftag{FOL}{\Sat{M(\Gamma^*)}}{\pSat{v(\Gamma^*)}}{\indfrm}$
    precies als zowel
    $\iftag{FOL}{\Sat{M(\Gamma^*)}}{\pSat{v(\Gamma^*)}}{!B}$ and
    $\iftag{FOL}{\Sat{M(\Gamma^*)}}{\pSat{v(\Gamma^*)}}{!C}$ gelden,
    volgens de definitie van vervulling. Volgens de inductieaanname is
    dat precies zo als zowel $!B \in \Gamma^*$ als $!C \in \Gamma^*$.
    Volgens \olref[ccs]{prop:ccs}\olref[ccs]{prop:ccs-and} geldt dat weer
    precies als $(!B \land !C) \in \Gamma^*$.}}}{}

\tagitem{prvOr}{%
\iftag{probOr}{%
  \indcase!{!A}{!B \lor !C}{}}{%
  \indcase{!A}{!B \lor !C}
    {$\iftag{FOL}{\Sat{M(\Gamma^*)}}{\pSat{v(\Gamma^*)}}{\indfrm}$
     precies als $\iftag{FOL}{\Sat{M(\Gamma^*)}}{\pSat{v(\Gamma^*)}}{!B}$ of
     $\iftag{FOL}{\Sat{M(\Gamma^*)}}{\pSat{v(\Gamma^*)}}{!C}$ geldt,
     volgens de definitie van vervulling. Volgens de inductieaanname is
     dat precies zo als $!B \in \Gamma^*$ of $!C \in \Gamma^*$.
     Volgens \olref[ccs]{prop:ccs}\olref[ccs]{prop:ccs-or} geldt dat weer
     precies als $(!B \lor !C) \in \Gamma^*$.}}}{}
  
\tagitem{prvIf}{%
\iftag{probIf}{%
  \indcase!{!A}{!B \lif !C}{}}{%
  \indcase{!A}{!B \lif !C}
    {$\iftag{FOL}{\Sat{M(\Gamma^*)}}{\pSat{v(\Gamma^*)}}{\indfrm}$
     precies als $\iftag{FOL}{\Sat/{M(\Gamma^*)}}{\pSat/{v(\Gamma^*)}}{!B}$
     of $\iftag{FOL}{\Sat{M (\Gamma^*)}}{\pSat{v(\Gamma^*)}}{!C}$ geldt,
     volgens de definitie van vervulling. Volgens de inductieaanname is
     dat precies zo als $!B \notin \Gamma^*$ of $!C \in \Gamma^*$.
     Volgens \olref[ccs]{prop:ccs}\olref[ccs]{prop:ccs-if} geldt dat weer
     precies als $(!B \lif !C) \in \Gamma^*$.}}}{}

\iftag{FOL}{%
\tagitem{prvAll}{%
  \iftag{probAll}{%
    \indcase!{!A}{\lforall[x][!B(x)]}{}}{%
    \indcase{!A}{\lforall[x][!B(x)]}{$\Sat{M(\Gamma^*)}{\indfrm}$ precies
      als $\Sat{M(\Gamma^*)}{!B(t)}$ voor elke term~$t$, volgens
      \olref{prop:quant-termmodel}. De inductieaanname zegt dat dit precies
      zo is als $!B(t) \in \Gamma^*$ voor elke term~$t$. Volgens
      \olref[hen]{prop:saturated-instances} geldt dat weer precies als
      $\lforall[x][!A(x)] \in \Gamma^*$.}}}{}

\tagitem{prvEx}{%
  \iftag{probEx}{%
    \indcase!{!A}{\lexists[x][!B(x)]}{}}{%
    \indcase{!A}{\lexists[x][!B(x)]}{$\Sat{M(\Gamma^*)}{\indfrm}$ precies
      als $\Sat{M(\Gamma^*)}{!B(t)}$ voor minstens één term~$t$, volgens
      \olref{prop:quant-termmodel}. De inductieaanname zegt dat dit precies
      zo is als $!B(t) \in \Gamma^*$ voor minstens één term~$t$. Volgens
      \olref[hen]{prop:saturated-instances} geldt dat weer precies als
      $\lexists[x][!B(x)] \in \Gamma^*$.}}}{}
}{}
\end{enumerate}
\end{proof}


\tagprob[FOL]{probOr,probAnd,probIf,probEx,probAll}

\begin{prob}
Maak het bewijs van \olref[fol][com][mod]{lem:truth} af.
\end{prob}

\tagendprob

\tagprob[notFOL]{probOr,probAnd,probIf}

\begin{prob}
Maak het bewijs van \olref[pl][com][mod]{lem:truth} af.
\end{prob}

\tagendprob

\end{document}
